\section{Operating Systems}
\label{sec:qna-os}

\subsection*{Processes and threads}

\question{Difference between a process and a thread?}
\answer{A process is an independent program in execution with its own
address space, memory, and resources. A thread is a lightweight unit of
execution \emph{within} a process; threads of one process share its address
space and resources but have their own stack and registers. Threads are
cheaper to create and switch, but a bug in one can corrupt the shared
space.}

\question{What is a context switch and why is it expensive?}
\answer{Saving the state (registers, program counter) of the currently
running process/thread and loading another's, so the CPU can switch
execution. It is costly because of the save/restore work plus cache and
TLB pollution -- the new task starts with cold caches.}

\question{What are the states in a process life cycle?}
\answer{New, Ready (waiting for CPU), Running (on CPU), Waiting/Blocked
(waiting for I/O or an event), and Terminated. The scheduler moves
processes between Ready and Running; I/O moves them to and from Waiting.}

\question{User mode vs kernel mode?}
\answer{Kernel mode has unrestricted access to hardware and all
instructions; user mode is restricted and cannot directly touch hardware.
Applications run in user mode and request privileged operations via
\emph{system calls}, which trap into the kernel. This protects the system
from buggy or malicious programs.}

\subsection*{Scheduling}

\question{Compare common CPU scheduling algorithms.}
\answer{\textbf{FCFS} -- simple, but long jobs delay others (convoy
effect). \textbf{SJF} -- optimal average waiting time but needs burst
prediction and can starve long jobs. \textbf{Round Robin} -- fair, good for
time-sharing, performance hinges on the time quantum. \textbf{Priority} --
important jobs first, risks starvation (fixed by aging).}

\question{What is starvation and how is it prevented?}
\answer{Starvation is a process waiting indefinitely because others are
continually preferred (e.g.\ low priority under priority scheduling).
\emph{Aging} prevents it by gradually raising a waiting process's priority
until it runs.}

\subsection*{Concurrency}

\question{What is a race condition?}
\answer{When two or more threads access shared data concurrently and the
outcome depends on the timing/interleaving of their operations, producing
incorrect results. Fixed by synchronising access to the \emph{critical
section} so only one thread modifies the shared data at a time.}

\question{Difference between a mutex and a semaphore?}
\answer{A mutex is a locking mechanism for mutual exclusion -- owned by one
thread, which must be the one to unlock it (binary). A semaphore is a
signalling counter allowing up to $N$ concurrent accesses; any thread may
signal. Use a mutex for exclusive access, a counting semaphore to limit a
resource pool.}

\question{What are the conditions for deadlock and how is it handled?}
\answer{Four must hold together: mutual exclusion, hold-and-wait, no
preemption, and circular wait. Handle by \emph{prevention} (break a
condition, e.g.\ lock ordering), \emph{avoidance} (Banker's algorithm),
\emph{detection and recovery}, or \emph{ignoring} it (the ``ostrich''
approach) if rare.}

\question{Difference between deadlock and livelock?}
\answer{In deadlock, processes are blocked and make no progress. In
livelock, they are active and keep changing state in response to each other
but still make no useful progress (like two people stepping aside in the
same direction repeatedly).}

\subsection*{Memory}

\question{What is virtual memory?}
\answer{An abstraction giving each process its own large, contiguous
address space independent of physical RAM, using paging and a disk-backed
swap area. It enables running programs larger than RAM, isolation between
processes, and simpler memory allocation.}

\question{What is paging, and what is a page fault?}
\answer{Paging divides memory into fixed-size pages (virtual) and frames
(physical), mapped by a page table, eliminating external fragmentation. A
\emph{page fault} occurs when a referenced page is not in RAM; the OS traps,
loads the page from disk (possibly evicting another), and resumes.}

\question{Difference between paging and segmentation?}
\answer{Paging splits memory into fixed-size pages (transparent to the
program, no external fragmentation but some internal). Segmentation splits
it into variable-size logical units (code, stack, heap) meaningful to the
program (external fragmentation possible). Modern systems often combine
both.}

\question{What is thrashing?}
\answer{When a system spends more time swapping pages in and out than
executing, because processes collectively need more memory than is
available. Detected via high paging activity; mitigated by reducing the
degree of multiprogramming or adding RAM (working-set model).}
